\section{A variance coordinate and its marginal decay}
\label{sec:variance}

The estimates of Proposition~\ref{mom:one-step} leave one quantity
uncontracted: the variance density $W$ of the energy charges.  We now
replace $W$ by a polynomial coordinate whose leading change under one
scale step is a negative multiple of $W^2$.  The normalization of a finite insertion
is responsible for that sign.

The coordinate's quadratic part will be exactly $W$.  In a product of
two independent centered insertions, the normalizing denominator first
affects the variance at degree four; its distant energy contribution
is negative.  Including the local cubic and quartic insertion
coefficients makes this calculation compatible with an exact scale
step.  We define those coefficients before computing their change.

Throughout this section the spacing is $s$, and $q_s=Q/s$.  All input
potentials are centered under the pure reference law $\E_0$, which
includes the auxiliary variables.  Environmental expectation $\E$ does
not integrate these auxiliary variables; it does integrate the independent
priorities and thresholds used to construct the scale map.  Write $R_0$
for a fixed dependence range in cell units, enlarged so that disjoint
sets of cells at distance greater than $R_0$ use disjoint environment
variables and have disjoint declared supports.

\subsection{Normalized insertion coefficients}

For a finite set of cell indices $I$ and real indeterminates
$\mathbf t=(t_i)_{i\in I}$, set
\begin{equation}
 F_I(\mathbf t)=\sum_{i\in I}t_i f_i,
 \qquad
 \mathcal T_I(\mathbf t)
 =q_s\left(\frac{e^{F_I(\mathbf t)}}{\E_0e^{F_I(\mathbf t)}}\right).
 \label{var:insertion}
\end{equation}
For each fixed environment these expressions are analytic near the
origin.  We use only their Taylor coefficients, so no exponential
integrability with respect to the environment is required.  In
particular, since $\E_0F_I=0$ and $q_s$ kills constants, the first three
homogeneous terms are
\begin{align}
 \mathcal T_1(F)&=q_s(F),&
 \mathcal T_2(F)&=\tfrac12q_s(F^2),\notag\\
 \mathcal T_3(F)&=\tfrac16q_s(F^3)
       -\tfrac12q_s(F)\E_0F^2.
 \label{var:charge-expansion}
\end{align}
Consequently the terms of degrees two, three, and four in the
environmental variance of $\mathcal T_I$ are, respectively,
\begin{equation}
 \Var(\mathcal T_1),\qquad
 2\Cov(\mathcal T_1,\mathcal T_2),\qquad
 \Var(\mathcal T_2)+2\Cov(\mathcal T_1,\mathcal T_3).
 \label{var:variance-expansion}
\end{equation}
All covariances in this section are environmental unless a subscript
$0$ is present.

Here is an explicit convention for the coefficients, including repeated
indices.  For an ordered tuple $\boldsymbol i=(i_1,\ldots,i_d)$, give
its $d$ occurrences separate indeterminates $t_1,\ldots,t_d$, even
when some $i_a$ coincide.  Define the symmetric multilinear coefficient
\begin{equation}
 C_d(\boldsymbol i)=
 \left.
 \partial_{t_1}\cdots\partial_{t_d}
 \Var\left[
 q_s\left(\frac{\exp(\sum_{a=1}^d t_a f_{i_a})}
                   {\E_0\exp(\sum_{a=1}^d t_a f_{i_a})}\right)
 \right]\right|_{\mathbf t=0},
 \qquad 2\le d\le4.
 \label{var:polarized-coefficient}
\end{equation}
The derivative in this formula means coefficient extraction using
\eqref{var:charge-expansion}--\eqref{var:variance-expansion}; thus it
is a finite sum of moments.  In particular its definition remains valid
when the random potentials have only the finite moments assumed in
Section~\ref{sec:moments}.

\begin{definition}[Variance coordinate]\label{var:coordinate}
For $H>R_0$, define
\begin{equation}
 v_H=\sum_{d=2}^4\frac1{d!}
 \sum_{i_2,\ldots,i_d\in\Z^2}
 \1_{\{\diam(0,i_2,\ldots,i_d)\le H\}}
 C_d(0,i_2,\ldots,i_d).
 \label{var:coordinate-formula}
\end{equation}
The diameter is Euclidean diameter in cell units.  The coordinate is
real and need not be nonnegative for an arbitrary field.
\end{definition}

The factorial and the ordered slots in
\eqref{var:coordinate-formula} are the usual Taylor convention.  An
equivalent definition sums over all ordered tuples in a large cell box,
divides by the number of cells, and takes the limit.  Stationarity and
the finite diameter cutoff give this equivalence.  This latter definition
also applies when an intermediate allocation is periodic on parent
cells: average over complete parent periods before passing to the limit.

\begin{lemma}\label{var:near-W}
For each fixed $H>R_0$, there is $C_H<\infty$, independent of $s$,
such that
\begin{equation}
 |v_H-W|\le C_H X^3\qquad\text{if }X\le1.
 \label{var:near-W-bound}
\end{equation}
\end{lemma}
\begin{proof}
The quadratic coefficient is
$C_2(i,j)=2\Cov(q_s(f_i),q_s(f_j))$.  Its contribution is therefore
exactly $W$, because covariances vanish past the dependence range.
Every remaining coefficient has three or four potential occurrences.
Their supports fit in a ball of radius $O(Hs)$.  The charge bound in
Theorem~\ref{pure:charge}, the boundedness of pure expectation, and
H\"older's inequality bound each coefficient of degree $d$ by
$C_H X^d$.  There are finitely many tuples with first slot zero and
diameter at most $H$.
\end{proof}

\subsection{The change of cutoff}

We first change $H$ to $LH$ without changing the field.  This calculation
contains the sign of the flow.  Later we compare this larger cutoff with
the coordinate of the output field.

\begin{lemma}[Shell calculation]\label{var:shell}
Fix finite constants $A,B$.  If $0<z\le1$, $X\le Az$, and
$Y\le Bz^2$, then
\begin{equation}
 v_{LH}-v_H
 =-2\kappa W^2\sum_{y\in\Z^2:\ H<|y|\le LH}|y|^{-2}
   +e_H z^4+O_H(z^5),
 \label{var:shell-formula}
\end{equation}
where $|e_H|$ is bounded by a deterministic function tending to zero
as $H\to\infty$.  The constants may depend on $A,B$ and the fixed
scale-step parameters, but are uniform in $s$.
\end{lemma}
\begin{proof}
Decompose $f_i=\bar f_i+\widetilde f_i$, where
$\bar f_i=\E f_i$.  An occurrence of $\widetilde f_i$ has size order
one, and an occurrence of $\bar f_i$ has size order two.  Products
of total size order at least five contribute $O_H(z^5)$ by
H\"older's inequality.  For the remaining coefficients connect two
occurrences when their cell indices are within $R_0$.  A connected
component containing at most four occurrences has bounded diameter.
Moreover, a component consisting of a single centered occurrence
has zero environmental expectation in every term of the covariance
expansion.  Thus, once $H$ is large, a nonzero coefficient in the shell
has precisely two distant groups: either a centered pair and one mean,
or two centered pairs.  Each has total size order four.  This argument
applies also to the products of means subtracted in a covariance.

Call the two groups $A$ and $B$.  Let $1+g_A$ and $1+g_B$ be their
separately normalized insertion densities, expanded formally in their
own indeterminates.  Every coefficient of $g_A,g_B$ is pure centered.
Define
\begin{equation}
 T_A=q_s(g_A),\quad T_B=q_s(g_B),\quad
 h_{AB}=q_s(g_Ag_B),\quad c_{AB}=\E_0(g_Ag_B).
 \label{var:two-group-quantities}
\end{equation}
Multiplying the two densities and normalizing again gives the exact
formal identity
\begin{equation}
 \mathcal T_{A\cup B}
       =\frac{T_A+T_B+h_{AB}}{1+c_{AB}}.
 \label{var:two-group-identity}
\end{equation}
If the anchor distance is $l$, Proposition~\ref{pure:fusion} gives,
coefficient by coefficient,
\begin{equation}
 h_{AB}=O(l^{-u}),\qquad
 c_{AB}=l^{-2}\bigl(\kappa T_AT_B+o(1)\bigr).
 \label{var:cross-coefficient-bounds}
\end{equation}
The estimates carry the product of the norms of the local coefficient
functions.  All group diameters are bounded, so their constants are
uniform in the slots and in $s$.

The variance of $T_A+T_B$ has no coefficient using both independent
groups.  Contributions with two cross factors $h_{AB}$ or $c_{AB}$
have shell sum bounded by $C H^{2-2u}z^4=o_H(1)z^4$.  For a single
$h_{AB}$ we use the stronger, symmetric part of
Proposition~\ref{pure:fusion}.  To justify its applicability, first
replace the cutoff on total diameter by the cutoff on anchor distance
and sum all bounded relative positions inside each group.  Environmental
factorization turns the required coefficient into a finite sum of
bilinear charge pairings of deterministic functions.  More explicitly,
the two possible local sums, besides $\bar f_x$, are
\begin{align*}
 D_x&=\sum_{|a|\le R_0}
       \E\bigl[q_s(\widetilde f_x)\widetilde f_{x+a}\bigr],\\
 M_x&=\sum_{|a|\le R_0}
       \E\bigl[\widetilde f_x\widetilde f_{x+a}
                   -\E_0(\widetilde f_x\widetilde f_{x+a})\bigr].
\end{align*}
The pair-plus-mean sector gives pairings $q_s(D_x\bar f_y)$;
the pair-plus-pair sector gives $q_s(D_xM_y)$, with their exchanges
and the fixed scalar factors from Taylor expansion.  This list follows
by taking the linear charge from one group and the term linear in that
group in $h_{AB}$; the other group supplies its mean or its quadratic
normalized density.  Terms leaving a centered singleton vanish as
above.  All displayed functions are pure centered, and $D_x,M_x$
have norm $O(z^2)$.  Rotation of the offset sum preserves both its
index set and the environment law, while charge is a rotation-invariant
scalar.  Thus each function is quarter-turn invariant about its anchor.
The resulting bound is
$C l^{-w}z^4$, whose shell sum is $O(H^{2-w})z^4=o_H(1)z^4$.

The replacement of diameter by anchor distance affects only strips
of bounded thickness at the two ends of the shell.  There are $O(H)$
possible anchor displacements in these strips, and the unsymmetrized
bound $O(H^{-u})$ gives an error $O(H^{1-u})z^4$.  Thus the replacement
and its reversal both have vanishing error as $H\to\infty$.

Only the single-$c_{AB}$ contribution remains.  Up to size order four
it is
\begin{equation}
 -2\Cov\bigl(T_A+T_B,(T_A+T_B)c_{AB}\bigr).
 \label{var:normalization-term}
\end{equation}
Its four occurrences must all be centered; a mean would increase
the size order.  Write $U=T_{A,1}$ and $V=T_{B,1}$.  These are
independent centered random scalars.  At bidegree $(2,2)$,
\begin{align}
 -2\Cov(U+V,(U+V)\kappa l^{-2}UV)
 &=-2\kappa l^{-2}\E[(U+V)^2UV]\notag\\
 &=-4\kappa l^{-2}\E U^2\E V^2.
 \label{var:quartic-sign}
\end{align}
The $o(1)$ in \eqref{var:cross-coefficient-bounds} gives
$o_H(1)z^4$ after shell summation, since $L$ is fixed and
$\sum_{H<|y|\le LH}|y|^{-2}$ is bounded.

We finally spell out the slot count.  Put
$a_{ij}=\Cov(q_s(f_i),q_s(f_j))$.  Suppose the four labeled slots
split into the far pairs $\{1,2\}$ and $\{3,4\}$.  The coefficient
of $t_1t_2$ in $\E U^2$ is $2a_{i_1i_2}$, even when the two cell
indices coincide; the occurrences still have distinct indeterminates.
The other pair contributes $2a_{i_3i_4}$.  Hence the polarized
normalization contribution is
\[
 -16\kappa |i_3-i_1|^{-2}a_{i_1i_2}a_{i_3i_4}.
\]
There are three partitions of four labeled slots into two unordered
pairs.  Their sums agree by permutation of the dummy slots, and a
tuple in a sufficiently distant shell has at most one such far
pairing.  The prefactor from \eqref{var:coordinate-formula} is
$1/4!$, so the numerical coefficient is
$-16\cdot3/4!=-2$.  For the pairing just displayed, anchoring
$i_1=0$ and writing $y=i_3$ gives
\[
 \sum_{i_2}a_{0i_2}\sum_{i_4}a_{y i_4}=W^2.
\]
Both sums are finite, and stationarity makes the second equal to $W$.
Changes of the cutoff caused by bounded shifts of anchors have the
strip bounds already proved.  A change of anchor in the kernel itself
also has vanishing error: for bounded $\delta$ and $|y|\asymp H$,
$\bigl||y+\delta|^{-2}-|y|^{-2}\bigr|\le C|\delta|H^{-3}$,
whose shell sum is $O(H^{-1})$.  This proves
\eqref{var:shell-formula}.
\end{proof}

\subsection{Passing to the output coordinate}

Let $f'_j$ be the output potentials at spacing $Ls$, and let $v'_H$
be their coordinate.  A finite insertion has the same exterior marginal
before and after the regular scale map.  The next argument explains why
this identity applies to the truncated, locally allocated coordinate.

\begin{lemma}[Compatibility of the coordinate]\label{var:compatibility}
Under the hypotheses of Lemma~\ref{var:shell},
\begin{equation}
 v'_H-v_{LH}=e_Hz^4+O_H(z^5),
 \label{var:compatibility-bound}
\end{equation}
with the same uniformity convention for $e_H$.
\end{lemma}
\begin{proof}
First use the regular map and retain the child indeterminates.  For
finitely many labeled inputs its normalized total insertion charge
obeys
\begin{equation}
 \mathcal T_{\mathrm{out}}=L^{-1}\mathcal T_{\mathrm{in}}
 \label{var:total-charge-preservation}
\end{equation}
by Corollary~\ref{map:insertion-charge}.  One can also see the normalization
directly: the conditional coin likelihood integrates to one at every
fixed old configuration, and each conditional integration preserves
pairing with a test outside the affected region.  Pairing with a distant
pure bond identifies $Q$ by Theorem~\ref{pure:charge}.  Centering
the allocated output logs changes only a scalar, which disappears from
the normalized insertion.  The factor $L^{-1}$ in
\eqref{var:total-charge-preservation} is exactly the change from $Q/s$
to $Q/(Ls)$.

This argument holds for every fixed realization of the fresh priorities.
Taking environmental variance and then Taylor coefficients therefore
preserves complete coefficients, including their mean subtractions.
The variance factor $L^{-2}$ cancels the $L^2$ children per parent
when sums are taken per unit area.  If an allocation distinguishes
positions within a parent, take the area average over complete parent
periods.  There is no assumption of invariance of individual allocated
coefficients under child translations.

We must compare the two cutoffs.  Every output Taylor coefficient uses
children in a fixed stencil.  Thus classifying a coefficient by parent
diameter $H$ rather than child diameter $LH$ changes its classification
only when the child diameter lies in a strip of width $O_L(1)$ about
$LH$.  Keep together the full covariance coefficient for each fixed
tuple of output arguments, including all its mean subtractions, and
only then expand its arguments in child slots.  The parent cutoff
is constant on precisely this grouping.  More explicitly, write
$R_j$ for the regular output at parent $j$, as a function of its
child potentials.  A Taylor contribution with $m$ output arguments
and child degree $d$, where $2\le m\le d\le4$, assigns its $d$
occurrence slots to nonempty blocks $B_1,\ldots,B_m$ and
uses the random functions
\begin{equation}
 G_a=D^{|B_a|}R_{j_a}(0)
        [f_{i_b}:b\in B_a],\qquad 1\le a\le m.
 \label{var:taylor-blocks}
\end{equation}
The directions in this derivative are supported at their indicated
child labels; repeated labels are allowed.  Multilinear substitution
gives the complete polarized covariance coefficient of
$G_1,\ldots,G_m$, with its Taylor factors and the single common
weight $\1_{\{\diam(j_1,\ldots,j_m)\le H\}}$.

Discard size orders at least
five, at cost $O_H(z^5)$, and connect child slots when their enlarged
stencils can share environment variables.  A singleton centered child
has zero expectation even conditional on
the fresh priorities: the Taylor kernel depends on those independent
priorities and takes the child function as a multilinear argument.
A singleton mean costs size order two.  Two components containing means
only cannot create a lower-order term through random allocation.  At the
only relevant order their output is the linear map of
Lemma~\ref{map:linearization}, which is independent of priorities;
its selection indicators are reference auxiliary variables, integrated
by $Q$.  In particular its charges are deterministic when its inputs
are deterministic.  We are therefore left with two distant components
of bounded diameter and size order two each.

Every output argument receiving a child slot lies in that slot's
bounded stencil.  These components consequently partition the relevant
output arguments into two distant groups as well: each block $B_a$
in \eqref{var:taylor-blocks} lies wholly in one component.  Apply
\eqref{var:two-group-identity} to the two finite lists of output
arguments $G_a$, with separate insertion indeterminates $u_a$.
The groups are
independent in the environment, including their fresh priorities.
For their separate charges the formal identity
\[
 \Var\bigl(T_A(\mathbf u_A)+T_B(\mathbf u_B)\bigr)
 =\Var T_A(\mathbf u_A)+\Var T_B(\mathbf u_B)
\]
has zero coefficient involving variables from both groups, regardless
of their separate means.  Thus the unmixed contribution cancels for
each fixed output tuple and each child-block partition, before any
sum over output labels.  Repeated child labels merely repeat derivative
directions inside a group; no within-group independence is used.
The remaining terms have a cross factor and hence
an $O_L(H^{-u})$ bound from Proposition~\ref{pure:fusion}.
The strip therefore costs
$O_L(H^{1-u})z^4=o_H(1)z^4$.  Keeping the covariance subtractions with
each coefficient is essential here and is permitted because the cutoff
depends only on the tuple of labels.

Finally apply Corollary~\ref{mom:taylor-substitution}.  For each fixed
output tuple, \eqref{var:variance-expansion} is a sum of bounded
multilinear expressions in at most four output potentials, including
products of separately taken expectations.  The charge bound on its
fixed-$H$ support makes the kernels uniformly bounded in $s$.  The
corollary applies to the finite union of all determining stencils,
including their status checks.  It therefore replaces the regular
Taylor output by the actual output and removes child degrees above
four with total error $O_H(X^5)=O_H(z^5)$.  In particular the highest
intermediate degree is $4\cdot4=16<P=40$, so the finite moment
assumption suffices.  This proves \eqref{var:compatibility-bound}.
\end{proof}

\begin{proposition}[Marginal variance flow]\label{var:flow}
Fix the scale-step parameters and finite constants $A,B$.  For every
sufficiently large fixed $H$, the polynomial coordinate of
Definition~\ref{var:coordinate} satisfies, whenever
$0<z\le1$, $X\le Az$, and $Y\le Bz^2$,
\begin{equation}
 v'_H-v_H
 =-2\kappa W^2\sum_{y\in\Z^2:\ H<|y|\le LH}|y|^{-2}
   +e_Hz^4+O_H(z^5),
 \label{var:flow-formula}
\end{equation}
where $|e_H|\to0$ as $H\to\infty$, uniformly in the lattice scale.
Moreover $|v_H-W|\le C_HX^3$ for $X\le1$.
\end{proposition}
\begin{proof}
Combine Lemmas~\ref{var:near-W}, \ref{var:shell}, and
\ref{var:compatibility}.
\end{proof}

The shell sum tends to $2\pi\log L$ as $H\to\infty$.  Thus the
leading coefficient in \eqref{var:flow-formula} is strictly negative.
In the next section the temperature is chosen so that the mean charge
stays small enough for this variance decay to continue at every scale.
